\section{Điều bài báo chưa giải quyết}
\label{sec:ch13-dieu-chua-giai-quyet}

Cái trần thu hẹp yêu cầu đặt lên một dụng cụ đo tương lai mà không thay thế được
dụng cụ ấy, và nó để hở một chỗ mà chính bài đặt tên. Ngoài hai điều đó, bài còn
để ngỏ bốn chỗ nữa, và ba trong bốn chỗ là do chính bài nêu.

\index{độ phức tạp Kolmogorov!không tính được}%
Chỗ thứ nhất là chỗ nặng nhất và bài chỉ nói nửa câu về nó. \tn{Độ phức tạp
Kolmogorov}{Kolmogorov complexity} không tính được, nên $\varepsilon_f$,
$\kappa_f$ và mọi chặn dựng trên chúng đều không tính được. Không ai cầm một mô
hình và một lời giải thích rồi kiểm tra xem cặp ấy nằm trong hay ngoài
\tn{miền khả thi}{feasibility region} của
hình~\ref{fig:ch13-mien-kha-thi}. Bài bắc cầu bằng
bảng~\ref{tab:ch13-do-phuc-tap-lop-mo-hinh}, thay $K$ bằng độ dài một cách mã
hóa cho từng lớp mô hình, và tự gọi lập luận ấy là phác thảo chứng
minh~\autocite{p26limits}. Cây cầu ấy chỉ chịu được tải theo một chiều: mỗi ô
trong bảng là một chặn trên cho $K$, nên nó dùng được ở chỗ nào cần chặn trên và
không dùng được ở chỗ nào cần chặn dưới. Đó đúng là chỗ hệ quả 3.8 sẩy chân, như
mục~\ref{sec:ch13-tran-cao-bao-nhieu} đã ghi, và cái sẩy chân ấy nằm ngay tại
điểm nối giữa lý thuyết và thực hành chứ không ở một góc phụ.

Chỗ thứ hai là chỗ chương này đã nói ngay từ đầu: bài không chạy thí nghiệm nào.
Bài tự xếp việc ấy vào công việc tương lai và còn nêu cụ thể thí nghiệm đầu tiên
nên chạy, là khớp các hàm đa thức đơn giản bằng cây quyết định với độ phức tạp
tăng dần rồi đối chiếu đường cong thu được với các chặn đã
chứng minh~\autocite{p26limits}. Một thí nghiệm như thế không kiểm chứng được
định lý, vì định lý đã đúng; cái nó kiểm chứng là những ước lượng của
bảng~\ref{tab:ch13-do-phuc-tap-lop-mo-hinh}, tức đúng cây cầu vừa nói.

Chỗ thứ ba là bước mô hình hóa mà mục~\ref{sec:ch13-ham-xap-xi} đã đánh dấu là
quan trọng nhất trong bài. Toàn bộ khái niệm \enquote{dễ diễn giải} được quy về
một con số bit, và không có gì trong bài hiệu chỉnh con số ấy. Không ai biết
$b_{\text{người}}$ bằng bao nhiêu, và câu hỏi khó hơn là liệu nó có phải một con
số hay không: người đọc không đọc chương trình, và hai lời giải thích cùng độ
dài mô tả có thể dễ hiểu khác hẳn nhau tùy chúng nói về cái gì.
Chương~\ref{ch:con-nguoi-vang-mat} đã cho thấy cảm nhận của người đổi cả dấu khi
đổi lĩnh vực, nên giả định rằng nó rút gọn được về một hằng số là một giả định
mạnh. Bài liệt kê việc xây những thước đo độ phức tạp bám sát quá trình nhận
thức của con người hơn vào phần nghiên cứu tương lai của mình, cùng với việc
khảo sát độ phức tạp tính toán của bài toán tìm lời giải thích tối ưu, việc đưa
phân phối đầu vào vào khung một cách tường minh, việc đưa vào các khái niệm giải
thích theo nhân quả và \tn{phản thực}{counterfactual}, và các mô hình giải thích
động hoặc tương tác~\autocite{p26limits}.

Chỗ thứ tư là quan sát của sách. Năm bài báo của danh mục mà Phần~IV đọc kỹ từ
chương~\ref{ch:do-do-trung-thuc} tới đây, tức các bài 21 tới 25 của thang đọc,
đều nêu cùng một thứ đang thiếu, một
dụng cụ đo độ trung thực đã được kiểm định, và nêu từ những hướng khác nhau: đo
thấy nó hỏng, ghi nhận rằng không có chuẩn dùng nó, gọi tên bước tiếp theo chưa
ai làm, đếm xem phép kiểm chứng với người hiếm tới đâu. Bài này không nêu nó
lần nào, và chỗ trống trong danh sách hạn chế của nó là một chỗ trống khác hẳn:
cái nó thiếu là một đại lượng tính được thay cho $K$ và một thí nghiệm.
Chương~\ref{ch:khoang-trong-nghien-cuu} lấy giao của các phát biểu hạn chế ấy,
nên một bài không giao với các bài kia cũng là một dữ kiện, và dữ kiện ấy nói
rằng cái trần nằm ngoài chuỗi lập luận mà phần còn lại của Phần~IV đang dựng.

Định lý~\ref{thm:ch13-khoang-cach-do-phuc-tap} còn để lại một chỗ trống mà
chương sau đi thẳng vào, và các bất đồng vô hình của
mục~\ref{sec:ch13-tran-noi-gi} là nửa thứ hai của chỗ trống ấy. Định lý khẳng
định rằng một lời giải thích đơn giản hơn mô hình phải lệch khỏi mô hình ở đâu
đó, và nó không nói ở đâu; với nó, mọi chỗ lệch là như nhau. Nhưng chỗ lệch
không như nhau với người đọc, và nhất là không như nhau với người được chọn nó.
Nếu khoảng cách độ phức tạp buộc phải có sai số, thì một bên có quyền quyết định
lời giải thích trông ra sao có thể dồn toàn bộ sai số ấy vào những chỗ không ai
kiểm, chia nhỏ nó tới mức không chỉ số nào dựa trên sai số đầu ra thấy được, rồi
trưng ra một lời giải thích sạch sẽ trên những chỗ có người nhìn.
Chương~\ref{ch:giai-thich-duoi-tan-cong} đọc các công trình đã dựng đúng phép
dồn ấy thành phương pháp tấn công.
